\section{战争和武装冲突}
\begin{enumerate}
  \item 区分战争和武装冲突。
        \begin{enumerate}
          \item 开始、结束均以意思表示为标志。\al{要出现“宣战”这个词；宣布对方为敌国 \yjt 武装冲突}。
          \item 开始。打不打不重要，单方宣相互宣，口头宣书面宣无所谓。
          \item 宣而不战是战争，不宣而战是非战争武装冲突。\al{俄乌冲突}。
        \end{enumerate}
  \item 战争开始的法律后果。外交、经济 2 个方面。
        \begin{enumerate}
          \item 领土条约持续有效；友好关系的政治性条约废止；其他条约暂停适用，战争结束自动恢复。
          \item 外交、领事关系当即断绝，特权与豁免截至离境。
          \item 经贸（包括民间）往来禁止。国家、商业、私人。
                \begin{enumerate}
                  \item 仅敌国在本国境内公产可没收，占领区仅可以征用。使馆、私人财产除外。
                  \item 海上遇到的敌国公、商船舶及其货物可以拿捕、没收。宗教、科考、医疗、慈善、探险的除外。
                \end{enumerate}
        \end{enumerate}

        \begin{kk}{战时中立国。战时中立不是永久中立。}
          \begin{enumerate}
            \item 不作为。不能直接参与战争、提供资金 \hz 武器等资助、不能介入。
            \item 防止。交战国利用中立国的领土从事有关的战争行为。
            \item 容忍。对你带来的不便。
          \end{enumerate}
        \end{kk}
  \item 作战规则。战争和武装冲突均适用。
        \begin{enumerate}
          \item 不分皂白。禁止不分军事目标、平民目标进行打击。
          \item 对武器的限制。安理会武器禁用清单、毒气生化、杀伤人员的地雷（\al{针对车辆 \dui}）、核武器。
          \item 禁止改变环境。引发地震、海啸、大规模水灾。
          \item 禁止背信弃义。选项出现“假装”字眼，但不禁止使用诈术。\al{假装受伤、投降 \cuo}。
          \item 平民、受难者、战俘。一个原则一个立即指导一切。人道主义原则。立即遣返，不得迟延。
        \end{enumerate}
  \item 国际刑事法院。
        \begin{enumerate}
          \item 惩罚战争罪犯（个人） \yjt 不管国家间争端 \yjt 自然不以争端双方同意为条件。\al{国际罪犯、国际犯罪 \cuo}。
          \item 常设、独立。除了 WTO 专家组，其他都常设。\al{是联合国所属的 \cuo}。
          \item 侵略、战争、灭绝、人类。侵略罪、战争罪、灭绝种族罪、危害人类罪。
          \item 管辖权依据。非常宽，沾边就管。
        \end{enumerate}
\end{enumerate}
